% Imporditud: tools/import_eksperimendid.py
% Allikas: Eesti füüsikaolümpiaadi eksperimendi ülesannete
%   kogu (Rael Kalda, 2023), ülesanne Ü165, lahendus L165.
\setAuthor{EFO žürii}
\setRound{lõppvoor}
\setYear{2005}
\setNumber{G E2}
\setDifficulty{5}
\setTopic{dünaamika}
\setVahendid{paberiga kaetud kaldpind, statiiv, alumiiniumist keha, mõõtejoonlaud, paberileht}
\setPunktid{14}
\setAllikas{Eksperimendi kogu Ü165, lahendus lk 124}
\setLahendusseis{toores}

\prob{Kaldpind}
Määrake soojushulk, mis eraldub keha algkiiruseta libisemisel üle paberiga kaetud kaldpinna, kui keha stardib kaldpinna ülemise serva juurest, mille kaldenurk on $\varphi$ = 27,5$^\circ$. Hinnake mõõteviga. Keha tihedus on $\rho$ = 2700 kg/m$^3$. Keha mõõtmed võite lugeda kaldpinna pikkusega võrreldes tühiselt väikeseks.
\textit{Märkus:} Töö teostamisel võtke arvesse, et hõõrdetegur võib sõltuda libisemiskiirusest. Õhu takistust pole tarvis arvestada.

\hint

\solu
Kaldpinda mööda alla libiseva keha esialgne potentsiaalne energia muundub keha kineetiliseks energiaks kaldpinna lõpus ja hõõrdejõu tööks Ep = Ek + A. Hõõrdejõu töö muundub soojuseks, mis eraldub teatud soojushulgana. Seega eraldunud soojushulga leidmiseks peame teadma keha kineetilist energiat kaldpinna lõpus. Kuna kaldpinna ulatuses $\mu$ = const, ei saa me kineetilise energia määramisel lähtuda klotsi libisemise keskmisest kiirusest, et sellest arvutada lõppkiirust v = 2vk. Lõppkiiruse leidmiseks asetame kaldpinna alumise otsa täpselt laua äärele ja laseme klotsil maha kukkuda. Maas olevale paberile kukkunud klots jätab paberile jälje, mille alusel saame määrata kukkumise koha kauguse lauast ja klotsi kiiruse kaldpinna lõpus. Tähistades laua kõrguse h, klotsi jälje kauguse lauast s, kiiruse vertikaalsuunalise komponendi vh ja kiiruse horisontaalsuunalise komponendi vs, saame vs = v cos $\alpha$ ja vh = v sin $\alpha$.

s = vst = vt cos $\alpha$, h = vht + gt2 = vt sin $\alpha$ + gt2 Asendades t saame avaldada $v_2$: . 2 2

$v_2$ = h cos2 $\alpha$ gs2 . $-$ s sin $\alpha$ cos $\alpha$

Eraldunud soojushulk on seega

mv2 Q = mgH $-$ .

H v

s

vv h

h

s

Klotsi massi saame, kui mõõdame klotsi küljed, arvutame klotsi ruumala ja kasutame massi leidmiseks tiheduse valemit m = $\rho$V . Arvatavasti valmistab raskusi kõrguse H mõõtmine, sest tuleb arvestada, et klotsi raskuskese asub kaldpinna otsast eemal ja ka kaldpinna paksust. Seega laua pinnast kõrgust mõõta ei tohi. Teostatud: klotsi külgede mõõtmine ja massi arvutamine, kõrguse mõõtmine, mõõtemääramatuse arvutamine, kaldpinna paigutamine õige nurga alla ja s mõõtmine.
\probend
